\section{Задание}

Поставьте задачу кластеризации, используя датасет из лабораторной работы с начала семестра\cite{labin2026pandas}: сформируем её от классификации. Деления выборки на \texttt{train}/\texttt{test} не будет --- решаем задачу без учителя. Попробуйте получить кластеры, которые соответствуют исходным классам, и задайте функционал качества, по которому будете определять эффективность (сравнивать состав кластеров и классов).

\begin{enumerate}[label=\textbf{Task \arabic*:}, leftmargin=*]

\item \textit{[1 балл]}

Выберите категориальный признак, который показывает классы объектов, пусть будет от трёх до шести классов, оставьте только объекты, которые относятся к ним. Далее оставляем только числовые признаки \textit{(часть тоже можете удалить, пусть далее используются 5--10 числовых признаков)}.

\item \textit{[3 балла]}

Реализуйте алгоритм \texttt{k-means} \textbf{собственной реализации} и примените его к подготовленным данным. Сравните полученные кластеры с исходными классами с помощью выбранного функционала качества.

\item \textit{[1 балл]}

Примените \texttt{k-means} \textbf{из библиотеки} \texttt{sklearn} и сравните результат с собственной реализацией и с исходными классами.

\item \textit{[1 балл]}

Примените любой вариант \textbf{иерархической кластеризации} (из любой библиотеки) и оцените соответствие кластеров исходным классам.

\item \textit{[1 балл]}

Примените \textbf{DBSCAN} (из любой библиотеки) и оцените соответствие кластеров исходным классам.

\item \textit{[1 балл]}

Попробуйте улучшить результат, используя \textbf{уменьшение размерности} (метод PCA --- метод главных компонент) из библиотеки как предобработку с одним из методов кластеризации.

\item \textit{[1 балл]}

Попробуйте улучшить результат, используя \textbf{нормализацию} (выберите подходящий вариант сами) как предобработку с одним из методов кластеризации.

\item \textit{[1 балл]}

Сделайте \textbf{выводы} по итогам работы: сравните методы между собой, оцените влияние предобработок (PCA и нормализации) и сделайте заключение о том, насколько хорошо кластеры соответствуют исходным классам.

\end{enumerate}